\section{Related Work}

\subsection{Surveys and position papers}

Existing surveys and position papers provide important context but typically emphasize either broad HPC storage evolution or computational storage offload, rather than GPU-centered storage-stack evolution as a unified architectural problem. For example, recent HPC storage surveys summarize filesystems, deployment models, and protocol trends, while computational-storage surveys focus on where computation should execute in the storage hierarchy; neither thread fully models GPU-driven control-path transitions across local and remote I/O domains \cite{zhang2025hpcstoragesurvey,traub2023compstoragesurvey,liu2017rackscale}. As a result, prior surveys are valuable baselines but leave a synthesis gap at the intersection of accelerator control, storage semantics, and workload-coupled data services \cite{murray2021tfdata,rajbhandari2021zeroinfinity}.

\subsection{Traditional CPU-orchestrated storage systems}

Traditional POSIX-oriented stacks and host-managed I/O software place the CPU on the critical path for request submission, metadata handling, and completion processing, which becomes increasingly expensive as accelerator throughput grows. Even when optimized user-space paths are used, control-plane work remains primarily host-centric, so GPUs still observe indirect storage service and coordination overheads \cite{linux_io_uring,spdk_overview,nvme_spec}. This CPU-first orchestration model is one root cause of low accelerator utilization under data-intensive pipelines, especially when many small operations amplify control and synchronization costs \cite{murray2021tfdata,hoard2018}.

\subsection{GPU-aware and GPU-integrated storage systems}

GPU-aware file and runtime systems moved beyond host-only orchestration by exposing storage abstractions to GPU execution contexts and by co-designing data access with accelerator behavior. GPUfs and follow-up work demonstrated that GPU-side file access is feasible but requires careful handling of consistency, granularity, and control interactions; newer systems such as GeminiFS further align filesystem behavior with GPU-dominated workflows \cite{silberstein2013gpufs,shahar2016gpufs,geminifs2025}. These efforts are architecturally significant because they begin to redefine storage as a subsystem co-evolving with accelerator scheduling, not merely a backend byte source \cite{qureshi2023bam,qiu2025geminifs}.

\subsection{GPU direct and bypass I/O mechanisms}

Bypass mechanisms (GPUDirect RDMA and GPUDirect Storage) substantially reduce host bounce buffering and improve transfer efficiency, and they are now foundational in production heterogeneous clusters. However, this line of work is primarily data-path-centric: it optimizes copy avoidance and transport efficiency but often assumes largely unchanged control and metadata planes above the transport layer \cite{nvidia_gdrdma,nvidia_gds_blog,nvidia_gds_overview,nvidia_gds_best_practices}. Consequently, bypass I/O is necessary but insufficient for explaining end-to-end architectural evolution from CPU-mediated orchestration to GPU-participatory storage control \cite{pcie_spec,infiniband_spec,nvmeof_spec}.

\subsection{Storage for AI workloads}

AI input pipelines, distributed caching, and model-state persistence reveal workload properties that differ sharply from assumptions in conventional enterprise and HPC storage tuning. tf.data and Hoard show that data loading, prefetching, and locality decisions can dominate training throughput, while memory-tiering/checkpoint work shows that persistence policies directly affect model scale and recovery behavior \cite{murray2021tfdata,hoard2018,rajbhandari2021zeroinfinity}. These results suggest that storage architecture for GPU clusters must be co-designed with pipeline structure and training dynamics, rather than optimized as an isolated block-I/O subsystem \cite{hdf5gds2020,geminifs2025}.

\subsection{Disaggregated and remote storage for accelerators}

Disaggregated storage and remote-access protocols provide elasticity and pooling that are increasingly relevant for accelerator clusters, where local capacity is insufficient for multi-tenant AI and HPC workloads. NVMe-oF/RDMA fabrics and rack-scale disaggregation studies establish the deployment substrate for remote SSD access, but GPU-visible semantics and control-path placement remain underdeveloped compared with transport maturity \cite{nvmeof_spec,infiniband_spec,liu2017rackscale}. Thus, remote storage research contributes critical building blocks while still leaving open questions about GPU-initiated scheduling, metadata locality, and end-to-end coordination \cite{qureshi2023bam,nvidia_gdrdma}.

\subsection{Summary and gap identification}

Across these lines of work, the dominant emphasis is on optimizing \emph{how data moves} (copy elimination, DMA efficiency, transport throughput), while comparatively fewer studies address \emph{how storage control evolves} when GPUs become active participants in I/O orchestration. In particular, control-plane evolution, GPU-resident storage participation, and full end-to-end co-design across data, metadata, and workload pipelines remain fragmented across separate communities and artifacts \cite{nvidia_gds_overview,silberstein2013gpufs,qureshi2023bam,murray2021tfdata}. Therefore, a unified survey is needed to connect CPU-orchestrated stacks, GPU-aware integration, direct/bypass paths, and emerging GPU-driven storage paradigms within one architectural framework \cite{geminifs2025,zhang2025hpcstoragesurvey,traub2023compstoragesurvey}.
